
\chapter{Альтернативные методы: QR, BNPL, цифровой рубль}~\label{ch:alt-methods}
QR, BNPL и цифровой рубль удобно сравнивать по трём
осям: кто подтверждает платёж, где возникает расчётная финальность и
на ком остаётся риск спора или невозвратности.

\section{Таксономия методов оплаты}~\label{sec:alt-taxonomy}

Когда продуктовая команда говорит \enquote{добавим альтернативный способ
оплаты}, она легко смешивает в одну корзину вещи разной природы:
карту, банковский перевод, кошелёк, рассрочку и цифровые деньги.
Таксономия нужна затем, чтобы сравнивать их по устройству, а не по
маркетинговому описанию.

Первый и исторически наиболее распространённый класс: карточные
платежи.
Здесь списание инициирует мерчант: он отправляет запрос по
карточным реквизитам клиента и получает либо авторизацию, либо отказ.
Сильные стороны этого класса: глобальная инфраструктура, chargeback
и зрелый контур PCI DSS; слабая сторона: стоимость. Сюда входят
Visa, Mastercard, Мир, UnionPay и AmEx.

Второй класс: переводы со счёта на счёт (account-to-account, A2A).
Здесь уже клиент подтверждает перевод со своего счёта на счёт мерчанта.
Примеры: СБП, SEPA Instant Credit Transfer, PIX и UPI~\cite{sbp-main,sepa-instant,pix,upi}.
Сильная сторона A2A: финальность и низкие комиссии; слабая сторона:
более узкая защита покупателя.

Третий класс: кошельки, то есть надстройка над базовым платёжным
контуром. Кошелёк замкнутого контура (closed-loop wallet) хранит деньги
внутри замкнутого контура, как у Alipay и WeChat Pay. Кошелёк
открытого контура (open-loop wallet) лишь оборачивает карточный
инструмент токеном, как Apple Pay или Google Pay.

Четвёртый класс: BNPL (Buy Now Pay Later). Провайдер оплачивает
мерчанту всю сумму сразу, а клиент возвращает её частями по
собственному графику. Снаружи этот
поток похож на эквайринг с иной комиссией, но его экономическая модель
другая: в центре стоит кредитное решение провайдера.

Пятый класс: наличный сценарий (cash-based). Платёж начинается
онлайн, но завершается офлайн через ваучер или платёжный референс.
Типичный пример: OXXO Pay в Мексике. Покупатель получает код
платежа и завершает операцию в физической точке~\cite{oxxo-pay}.
Главный параметр здесь: задержка подтверждения. Деньги не приходят
мгновенно.

Шестой класс: цифровые деньги вне банковского депозита. Внутри этого
класса выделяются по меньшей мере две подгруппы. Цифровая валюта
центрального
банка (central bank digital currency, CBDC) означает цифровую форму
обязательства центрального банка. Криптовалюты и стейблкоины, напротив,
строятся на иной модели эмиссии и регулирования. Внешне их иногда
ставят в один ряд только потому, что и те, и другие не сводятся
к обычному банковскому счёту.

\begin{table}[tbp]
  \begingroup
  \scriptsize
  \setlength{\tabcolsep}{2pt}
  \renewcommand{\arraystretch}{1.2}
  \makebox[\textwidth][l]{%
  \begin{minipage}{\fullwidth}
  \begin{tabularx}{\fullwidth}{
      @{}>{\RaggedRight\arraybackslash\bfseries}p{0.170\fullwidth}
      >{\RaggedRight\arraybackslash}p{0.075\fullwidth}
      >{\RaggedRight\arraybackslash}p{0.095\fullwidth}
      >{\RaggedRight\arraybackslash}p{0.090\fullwidth}
      >{\RaggedRight\arraybackslash}p{0.085\fullwidth}
      >{\RaggedRight\arraybackslash}p{0.145\fullwidth}
      >{\RaggedRight\arraybackslash}p{0.100\fullwidth}
      >{\RaggedRight\arraybackslash}p{0.155\fullwidth}@{}
    }
    \toprule
    \rowcolor{Panel}
    Метод оплаты & Тип & Финальность & Chargeback & PCI
    scope & Международность & Скорость & Комиссия мерчанта \\
    \midrule
    \rowcolor{Soft}
    Карты &
    \cellcolor{Bad!10}pull &
    \cellcolor{Bad!10}нет &
    \cellcolor{Bad!10}есть &
    \cellcolor{Bad!10}да &
    \cellcolor{Good!14}глобальный &
    \cellcolor{Good!14}мс &
    \cellcolor{Bad!10}высокая \\
    A2A (СБП/PIX) &
    \cellcolor{Good!14}push &
    \cellcolor{Good!14}да &
    \cellcolor{Good!14}нет &
    \cellcolor{Good!14}нет &
    \cellcolor{Ink!6}региональный &
    \cellcolor{Good!14}сек &
    \cellcolor{Good!14}низкая \\
    \rowcolor{Soft}
    Closed-loop кошелёк &
    \cellcolor{Ink!6}push/pull &
    \cellcolor{Ink!6}зависит &
    \cellcolor{Ink!6}ограничен &
    \cellcolor{Good!14}нет &
    \cellcolor{Bad!10}нишевый &
    \cellcolor{Good!14}мс &
    \cellcolor{Ink!6}средняя \\
    BNPL &
    \cellcolor{Bad!10}pull &
    \cellcolor{Bad!10}нет &
    \cellcolor{Bad!10}есть &
    \cellcolor{Ink!6}нет &
    \cellcolor{Ink!6}региональный &
    \cellcolor{Good!14}сек &
    \cellcolor{Ink!6}средняя \\
    \rowcolor{Soft}
    Cash-based &
    \cellcolor{Good!14}push &
    \cellcolor{Good!14}да &
    \cellcolor{Good!14}нет &
    \cellcolor{Good!14}нет &
    \cellcolor{Bad!10}нишевый &
    \cellcolor{Bad!10}часы / дни &
    \cellcolor{Ink!6}нет данных \\
    CBDC (цифровой рубль) &
    \cellcolor{Good!14}push &
    \cellcolor{Good!14}да &
    \cellcolor{Good!14}нет &
    \cellcolor{Good!14}нет &
    \cellcolor{Bad!10}региональный &
    \cellcolor{Ink!6}сек / мин &
    \cellcolor{Ink!6}нет данных \\
    \rowcolor{Soft}
    Crypto &
    \cellcolor{Good!14}push &
    \cellcolor{Good!14}да &
    \cellcolor{Good!14}нет &
    \cellcolor{Good!14}нет &
    \cellcolor{Ink!6}глобальный / нишевый &
    \cellcolor{Ink!6}мин / часы &
    \cellcolor{Ink!6}нет данных \\
    \bottomrule
  \end{tabularx}
  \par\smallskip
  \fcolorbox{Ink!18}{Soft}{%
    \parbox{\dimexpr\fullwidth-2\fboxsep-2\fboxrule\relax}{%
      \textcolor{Muted}{Зелёный цвет показывает сравнительное
        преимущество, красный указывает на ограничение, серый на нейтральную или
        зависящую от провайдера зону.
        * Данные оценочные и меняются в зависимости от конкретного
      провайдера, маршрута и географии.}%
    }%
  }
  \end{minipage}%
  }
  \endgroup
  \caption{\enquote{Методы оплаты / параметры}}~\label{fig:alt-taxonomy-matrix}
\end{table}

Вопросы \enquote{кто инициирует платёж}, \enquote{где наступает финальность} и \enquote{кто
несёт риск} переводят маркетинговые ярлыки в сравнимые архитектуры.

\section{QR-платежи: три модели}~\label{sec:alt-qr}

Фраза \enquote{мы принимаем QR} сообщает меньше, чем кажется. Пока не
назван платёжный контур, из неё неясно ни кто подтверждает платёж,
ни где наступает финальность, ни кто потом разбирает спорные операции.
QR-код сам по себе -- лишь способ передать платёжные данные без
NFC-оборудования или ручного ввода реквизитов. За ним может стоять
A2A-перевод, кошелёк замкнутого контура или карточная транзакция.

Физически QR-код в платежах используется в двух режимах.
Merchant-presented mode (MPM): мерчант показывает код, клиент
сканирует его камерой телефона. Customer-presented mode (CPM):
клиент открывает динамический код в приложении, а мерчант считывает
его сканером на кассе. MPM дешевле в развёртывании -- достаточно
напечатанной наклейки или экрана на кассе. CPM быстрее на высоком
потоке покупателей, но требует от мерчанта считывающее
оборудование. Обе модели существуют внутри каждого из трёх контуров
ниже, однако пропорция их использования различается.

\subsection*{СБП: QR поверх A2A-перевода}

В СБП QR-код ведёт в мгновенный перевод со счёта клиента на счёт
мерчанта через инфраструктуру НСПК~\cite{sbp-main}. Контур работает
как push-платёж: клиент инициирует и подтверждает списание сам.

Статический QR -- это напечатанный код, который содержит
идентификатор мерчанта, но не содержит сумму. Клиент сканирует код
в приложении своего банка, вводит сумму вручную и подтверждает перевод.
Такой формат годится для малого бизнеса: не нужно кассовое ПО,
достаточно наклейки. Динамический QR генерируется кассовой системой
под конкретную покупку и содержит реквизиты получателя, сумму
и назначение платежа. Клиент сканирует код, видит предзаполненные
данные и подтверждает одним действием. Кассовая система получает
обратный вызов об успешном переводе за несколько
секунд~\cite{sbp-faq-business}.

Помимо QR, СБП поддерживает оплату по платёжной ссылке, через NFC
и через привязку счёта в приложении
СБПэй~\cite{sbpay-app,sbp-main}. Все эти каналы ведут к одному
и тому же A2A-переводу внутри инфраструктуры НСПК; различается
только способ, которым клиент передаёт согласие на списание.

Тарифы С2В-переводов в СБП устанавливает Совет директоров Банка
России. Ставка межбанковского вознаграждения зависит от MCC мерчанта
и на практике составляет от 0{,}4\% до 0{,}7\% от суммы
перевода~\cite{cbr-sbp-c2b-tariffs}. Для сравнения: interchange
в карточном эквайринге обычно начинается от 1\% и может превышать 2\%.
Регулируемая ставка -- осознанное решение Банка России, направленное на
то, чтобы снизить стоимость приёма платежей для бизнеса.

Механизма chargeback в СБП нет: перевод финален с момента зачисления
на счёт получателя. Если клиент хочет вернуть деньги за покупку,
мерчант инициирует обратный перевод через свой
банк~\cite{sbp-faq-business}. Ошибочный перевод (например, не тому
получателю) также требует согласия получателя на возврат, если иной
порядок не установлен законом или решением
суда~\cite{cbr-sbp-faq-mistaken-transfer}. Эта модель выгоднее
мерчанту, но смещает риск спора на покупателя -- прямая
противоположность карточному chargeback.

\subsection*{Alipay и WeChat Pay: QR поверх замкнутого
кошелька}~\label{sec:alt-qr-wallet}

В Alipay и WeChat Pay QR-код ведёт не в банковский перевод,
а в платёж внутри замкнутого контура кошелька. Средства пользователя
хранятся на балансе кошелька (пополняемом с банковского счёта или
карты), и списание происходит на стороне оператора кошелька без
выхода в межбанковский клиринг.

Оба режима -- MPM и CPM -- здесь активно используются.
В документации Alipay+ merchant-presented mode описан как основной
сценарий для сторонних мерчантов: мерчант показывает QR-код,
клиент сканирует его в приложении кошелька, оператор списывает
средства с баланса клиента и зачисляет их на счёт
мерчанта~\cite{alipayplus-mpm}. Customer-presented mode распространён
в высокопоточных офлайн-точках (сетевые магазины, автоматы,
общественный транспорт): клиент показывает динамический штрихкод
из приложения, кассир сканирует его, и оплата проходит без
дополнительного подтверждения со стороны клиента при сумме ниже
установленного оператором порога.

Оператор кошелька совмещает роли, которые в карточном контуре
разделены между эмитентом и эквайрером. Он ведёт балансы и клиентов,
и мерчантов, проводит расчёт внутри своего реестра и сам определяет
правила диспутов. Для клиента это означает, что весь платёжный
опыт -- от пополнения до возврата -- замкнут внутри одного
приложения. Для мерчанта -- что интеграция идёт через API оператора
кошелька, а не через банк-эквайрер.

Трансграничное расширение работает через партнёрскую сеть. Alipay+
объединяет локальные кошельки (GCash, TrueMoney, KakaoPay и другие),
позволяя их пользователям платить в точках, принимающих Alipay, за
пределами домашней страны. Мерчант подключается к одному контракту
с локальным партнёром Alipay+ и получает доступ к пользователям
нескольких кошельков~\cite{alipayplus-mpm}. Такая модель отличается
от карточной глобальности Visa или Mastercard: вместо единого
глобального протокола -- федерация кошельков, каждый из которых
работает на собственных рельсах.

\subsection*{PIX: QR, встроенный в национальную платёжную
инфраструктуру}~\label{sec:alt-qr-pix}

PIX -- система мгновенных платежей Banco Central do Brasil,
запущенная в ноябре 2020 года~\cite{pix}. Как и СБП, PIX строится
вокруг мгновенного A2A-перевода с финальностью в секундах и низкой
комиссией. Расчёт происходит в Sistema de Pagamentos Instantâneos
(SPI) -- выделенной инфраструктуре мгновенных платежей,
управляемой центральным банком~\cite{pix-spi}.

Участие обязательно для всех финансовых институтов с более чем 500
тыс.\ активных счетов. Это отличает PIX от европейских A2A-систем,
где подключение к мгновенным платежам долгое время оставалось
добровольным: обязательность дала PIX покрытие с момента запуска.

QR-код в PIX -- штатный интерфейс самой системы, а не надстройка
PSP или банка. Banco Central do Brasil определяет два типа QR:
статический (многоразовый, содержит идентификатор получателя
и опционально сумму) и динамический (одноразовый, генерируется под
конкретную транзакцию с полным набором данных). Помимо сканирования
QR, пользователь может воспользоваться режимом Pix Copia
e Cola -- текстовым представлением тех же данных, которое копируется
и вставляется в банковское приложение~\cite{pix-faq}. Этот режим
решает сценарий, в котором и платёж, и QR-код находятся на одном
устройстве и физическое сканирование камерой невозможно.

Комиссия для физических лиц за переводы и платежи через PIX
равна нулю -- это установлено регуляторно. Для юридических лиц
тарифы определяются банком-участником, но на практике значительно
ниже карточного эквайринга. Как и в СБП, chargeback отсутствует:
PIX работает как push-платёж, и перевод финален с момента зачисления
на счёт получателя.

\subsection*{Сравнение трёх моделей}

\begin{table}[tbp]
  \begingroup
  \scriptsize
  \setlength{\tabcolsep}{3pt}
  \renewcommand{\arraystretch}{1.2}
  \makebox[\textwidth][l]{%
  \begin{minipage}{\fullwidth}
  \begin{tabularx}{\fullwidth}{
      @{}>{\RaggedRight\arraybackslash\bfseries}p{0.16\fullwidth}
      >{\RaggedRight\arraybackslash}p{0.27\fullwidth}
      >{\RaggedRight\arraybackslash}p{0.27\fullwidth}
      >{\RaggedRight\arraybackslash}p{0.27\fullwidth}@{}
    }
    \toprule
    \rowcolor{Panel}
    & СБП & Alipay / WeChat Pay & PIX \\
    \midrule
    \rowcolor{Soft}
    Контур & A2A-перевод (межбанковский) & Замкнутый кошелёк & A2A-перевод (межбанковский) \\
    Оператор & НСПК (Банк России) & Ant Group / Tencent & Banco Central do Brasil \\
    \rowcolor{Soft}
    Расчёт & Через НСПК, секунды & Внутри реестра оператора & SPI при BCB, секунды \\
    Финальность & Мгновенная (push) & Мгновенная (внутри контура) & Мгновенная (push) \\
    \rowcolor{Soft}
    Chargeback & Нет & Правила оператора & Нет \\
    Статический QR & Да & Да & Да \\
    \rowcolor{Soft}
    Динамический QR & Да & Да & Да \\
    Основной QR-режим & MPM & MPM + CPM & MPM \\
    \rowcolor{Soft}
    Комиссия С2B & 0{,}4--0{,}7\% (регулятор) & Устанавливает оператор & 0\% (физ.\ лица) \\
    Трансграничность & Региональный & Федерация кошельков & Региональный \\
    \bottomrule
  \end{tabularx}
  \end{minipage}%
  }
  \endgroup
  \caption{Три модели QR-платежей}~\label{tab:qr-three-models}
\end{table}

Три модели используют один и тот же физический канал -- QR-код, --
но за ним стоят разные контуры с разными свойствами (табл.~\ref{tab:qr-three-models}).
СБП и PIX проводят платёж через межбанковскую инфраструктуру, управляемую
центральным банком; кошелёк Alipay или WeChat Pay замыкает платёж
внутри собственного реестра. Первая архитектура даёт мгновенный
push-перевод с регуляторно-ограниченной комиссией и отсутствием
chargeback. Вторая даёт контроль оператора над всем циклом платежа,
включая правила диспутов, ценообразование и трансграничное расширение.

QR -- лишь упаковка для платёжного контура. Продуктовый вопрос
не в том, какой QR показать, а в том, какой контур стоит за этим QR
и какие свойства он приносит: финальность, стоимость, модель диспутов
и географию покрытия.

\section{BNPL на уровне протокола}~\label{sec:alt-bnpl}

BNPL (Buy Now Pay Later) со стороны клиента выглядит как рассрочка:
выбрал товар, разбил оплату на части, забрал покупку сразу. Но на
уровне протокола перед нами не разновидность карты и не банковский
перевод, а точка входа в кредитный продукт. В момент оплаты
BNPL-провайдер принимает кредитное решение в реальном времени:
оценивает клиента, одобряет или отклоняет заявку, и в случае
одобрения перечисляет мерчанту полную сумму покупки сразу. Клиент
затем возвращает деньги провайдеру равными частями по графику.

Мерчант получает выручку так же быстро, как при обычном эквайринге,
но платит за это комиссию BNPL-провайдеру. Кредитный риск целиком
лежит на провайдере: если клиент не выплатит оставшиеся части,
мерчант денег не потеряет. Это распределение рисков и определяет
экономику BNPL: провайдер зарабатывает на комиссии мерчанта
и иногда на пенях за просрочку, а мерчант получает рост конверсии
и увеличение среднего чека за счёт более низкого порога входа
для покупателя.

На уровне протокола существуют два паттерна.

\subsection*{BNPL-как-PSP}

В этой модели BNPL-провайдер работает как отдельный PSP. Мерчант
заключает с ним merchant agreement, встраивает его SDK или
перенаправление в чекаут и обрабатывает обратные вызовы с результатом
кредитного решения.

Поток выглядит так. Клиент на странице оплаты выбирает
BNPL-провайдера. Провайдер показывает клиенту условия: количество
платежей, даты, сумму каждого платежа. Клиент подтверждает. Провайдер
проводит скоринг и выносит решение в реальном времени. При одобрении
провайдер перечисляет мерчанту полную сумму покупки за вычетом комиссии.
Мерчант видит одну транзакцию на полную сумму и отгружает товар.
Далее провайдер самостоятельно списывает с клиента оставшиеся
платежи по графику -- мерчант в этом процессе не участвует.

Интеграция похожа на подключение обычного PSP: API
для создания заказа, серверные уведомления (webhook) об изменении
статуса, API для возвратов. Спорные сценарии определяются договором
с провайдером, а не правилами платёжной системы. Карточного
chargeback здесь нет, потому что карточная сеть в контуре
не участвует.

По этой модели работают Klarna, Affirm, Tabby и все основные
российские BNPL-сервисы.

\subsection*{BNPL-на-карте}

Второй паттерн встраивает рассрочку в карточный контур. Здесь
клиент расплачивается картой, а разбиение на части происходит
на стороне эмитента или отдельного кредитора внутри карточной сети.
Мерчант может вообще не знать, что покупка оплачена в рассрочку.

Visa и Mastercard реализуют это по-разному. Visa Installments
строится на MIT-фреймворке (Merchant-Initiated Transaction).
Первая авторизация проходит как обычная карточная операция
с хранимыми данными (stored credential) и POS Environment
кодом~I (installment MIT). После неё мерчант или
планировщик Visa последовательно выставляет транши по сохранённым
реквизитам. Кредитный риск лежит на эмитенте: именно он решает,
для каких карт и каких покупок доступна рассрочка. Мерчант должен
поддерживать Installments API у своего процессора, а эмитент должен
активировать программу для своих карт -- без этого предложение
рассрочки клиенту не появится~\cite{visa-installments}.

Mastercard Installments в основной реализации работает иначе.
Участвующий кредитор -- банк или BNPL-финтех, контрактованный
с Mastercard -- выпускает одноразовый виртуальный номер карты на
отдельном BIN. Клиент расплачивается этим номером у мерчанта.
Мерчант получает полную сумму сразу как обычную карточную
транзакцию. Кредитор самостоятельно собирает взносы
с клиента~\cite{mc-installments}. Мерчанту не нужно
подключать отдельный API или активировать программу, однако
по выделенному BIN транзакцию можно идентифицировать.

Различие между двумя паттернами меняет интеграционные требования.
BNPL-как-PSP требует отдельного merchant agreement с провайдером:
нельзя просто добавить его как метод оплаты без нового договора.
BNPL-на-карте требует поддержки Installments API у процессора
(Visa) или наличия участвующего кредитора в сети (Mastercard),
но не требует от мерчанта отдельной интеграции с BNPL-провайдером.

\subsection*{Российский BNPL}

Российский рынок BNPL сформировался по модели BNPL-как-PSP.
Отличие от западных аналогов в том, что все крупные российские
BNPL-сервисы принадлежат банкам или финансовым группам, у которых
уже есть банковская лицензия и кредитная инфраструктура. При этом
сервисы принимают карты любых банков, а не только собственных
клиентов -- ограничивающим фактором является скоринг, а не
эмитент карты.

\enquote{Долями} (Т-Банк) -- первый российский BNPL-сервис,
запущенный в апреле 2021 года~\cite{tbank-dolyame}. Мерчант
интегрирует его как внешний PSP: отдельный договор, API или
плагин для CMS, отдельный поток
уведомлений~\cite{tbank-dolyame-business}. \enquote{Подели}
устроен аналогично; формальный оператор сервиса -- ООО
\enquote{А-Четыре Технологии}, 100\%-дочернее общество
АО \enquote{Альфа-Банк}; интеграция через API, плагины для CMS
или через платёжных
партнёров~\cite{alfa-podeli,alfa-podeli-business}.
\enquote{Плати частями} (Сбер) -- отдельный BNPL-сервис,
оператором которого является ООО \enquote{Центр новых финансовых
сервисов} (ЦНФС), 100\%-дочернее общество Сбербанка. Мерчанту
необходимо подключить его как самостоятельный платёжный способ
у процессора -- он не связан с интеграцией
SberPay~\cite{sber-pay-parts}.

Перечисленные сервисы объединяет одна архитектура: провайдер
работает как PSP с собственным merchant agreement и скорингом,
а мерчант получает полную сумму покупки за вычетом комиссии.
Различие -- в модели подключения. \enquote{Долями}
и \enquote{Подели} требуют отдельного договора с оператором.
\enquote{Сплит} (Яндекс) работает внутри Yandex Pay SDK:
мерчант подключает его через единую интеграцию с Yandex Pay,
не заключая отдельного BNPL-договора~\cite{yandex-split}.
\enquote{Плати частями} подключается как отдельный метод
у процессора.

\subsection*{Регулирование}

До недавнего времени BNPL в большинстве юрисдикций существовал
в регуляторном пробеле: формально это не кредит (нет процентов
для клиента), но по экономической сути -- кредитный продукт.

В России этот пробел закрыт Федеральным законом от 31 июля
2025 года 283-ФЗ~\cite{fz-283}, который вступает в силу
1 апреля 2026 года. Закон вводит понятие \enquote{оператор
сервиса рассрочки} (ОСР) и устанавливает для него отдельный
регуляторный режим: регистрация в реестре Банка
России~\cite{cbr-osr}, требования к капиталу, обязательное
раскрытие полной стоимости рассрочки для клиента, ограничение
максимального срока и запрет на взимание вознаграждения
с клиента. При совокупной задолженности клиента свыше
50~тыс.\ рублей данные передаются в бюро кредитных историй.
Провайдер, который хочет продолжать работу после 1 апреля
2026 года, должен получить статус ОСР.

В ЕС BNPL включён в обновлённую Consumer Credit
Directive~\cite{eu-ccd-2023}: обязательная оценка
кредитоспособности, единый стандарт раскрытия условий, право
клиента на отказ.

Для продуктовой команды регуляторный сдвиг означает конкретную
вещь: провайдер, не получивший статус ОСР или не прошедший
европейскую авторизацию, перестанет работать. При выборе
BNPL-провайдера стоит проверить его регуляторный статус на дату
интеграции.

\practicenote{%
При выборе BNPL-провайдера проверьте три параметра:
долю одобрений среди вашей аудитории,
модель ответственности за диспуты и регуляторный статус провайдера
(для России -- наличие в реестре ОСР).

}
\section{Цифровой рубль: архитектура и сценарии}~\label{sec:alt-digital-ruble}

Цифровой рубль является цифровой формой национальной валюты, эмитированной
Банком России непосредственно~\cite{fz-340,cbr-digital-ruble},
то есть третьей формой денег наряду с наличными и безналичными
(см.\ раздел~\ref{sec:banks-forms-of-money}),
с отдельным реестром обязательств на платформе Банка России.
Сравнение с СБП подробно разобрано в разделе~\ref{sec:sbp-cbdc}.

В действующей розничной модели архитектура цифрового рубля
двухуровневая. Банк России эмитирует цифровой рубль и управляет
платформой, а банки дают клиентам доступ к кошелькам и исполняют
их поручения на платформе Банка России. Клиент управляет кошельком
через мобильное приложение банка, подключённого к платформе~\cite{cbr-digital-ruble}.
Требования к защите информации для банков, участников платформы,
установлены Положением Банка России
833-П~\cite{cbr-833p}.

В чём отличие от безналичного рубля? В природе обязательства. Безналичный
рубль -- обязательство коммерческого банка перед клиентом; цифровой
рубль -- прямое обязательство Банка России. Отсюда и вывод: для клиента
исчезает кредитный риск банка, а для банков возникает дискуссия о том,
как такая форма денег повлияет на депозитную базу~\cite{cbr-digital-ruble}.

Смарт-контракты реализованы в двух режимах с разной степенью готовности.
Простые смарт-контракты уже работают в пилотном контуре: они позволяют
автоматически переводить цифровые рубли в заданные дату и время.
К маю 2025 года в пилоте исполнено более 17~тыс.\ таких контрактов~\cite{cbr-digital-ruble}; типовые сценарии: автоплатёж за
жилищно-коммунальные услуги, автопополнение кошелька, поэтапная оплата
подрядчику~\cite{gazprombank-smart-contracts-2025}.

Коммерческая платформа смарт-контрактов, на которой банки смогут
публиковать собственные сценарии с конфиденциальными условиями и
юридической обёрткой, находится в стадии разработки. В марте 2026 года
глава Банка России сообщила, что регулятор \enquote{представит для
обсуждения концепцию платформы коммерческих смарт-контрактов} в первом
полугодии 2026 года~\cite{comnews-nabullina-smart-contracts-2026}.
Потенциальные применения: целевые выплаты, страховые сценарии, B2B
расчёты с условиями. Пока они описаны на уровне концепции. Офлайн-режим
также остаётся перспективной возможностью, а не запущенным функционалом~\cite{cbr-digital-ruble}.

Цифровой рубль нередко ставят рядом с СБП, потому что оба сюжета
связаны с публичной инфраструктурой и оба обещают более дешёвый
платёж. Но на архитектурном уровне это разные вещи. СБП переводит
безналичные рубли между банковскими счетами, тогда как цифровой рубль
живёт в отдельном кошельке как самостоятельная форма денег. Смешивать их
в одну модель нельзя~\cite{cbr-digital-ruble}.

\section{Почему карточные контуры пока доминируют}~\label{sec:alt-card-dominance}

Карты сохраняют центральное место в потребительских платежах из-за
наложения нескольких инерций.

Первая причина: сетевые эффекты. Глобальные карточные схемы уже
выстроили массовую эмиссию, приём и привычный пользовательский сценарий.
Мерчант, который принимает карты, автоматически доступен большинству
покупателей. Новый метод должен одновременно получить поддержку
со стороны эквайера и эмитента, а это самый трудный вариант сетевого
эффекта.

Вторая причина: международность. Карта, выпущенная в одной стране,
обычно работает во множестве других по единому протоколу. Локальные
A2A-системы, такие как СБП, PIX и UPI, обычно не взаимозаменяемы.

Третья причина: защита покупателя. Chargeback даёт механизм, который
A2A-системы не воспроизводят без изменения самой идеи финальности.
Покупатель, оплативший картой незнакомому мерчанту, имеет формальный
возвратный контур; покупатель, оплативший через A2A, чаще остаётся с
договорной защитой без схемного chargeback.

Наконец, четвёртая причина: инфраструктурная инерция. Мировая сеть
POS-терминалов уже выстроена под карточные схемы. Переоснащение
или добавление нового метода окупается только при массовом спросе, а
массовый спрос не возникает, пока мерчанты не видят достаточного
пользовательского потока. Это классический замкнутый круг внедрения.

Альтернативные методы выигрывают в нишах: низкая комиссия для
A2A, доступность для небанковских клиентов в наличных сценариях,
условные платежи в контуре CBDC, рассрочка в BNPL и программируемые
трансграничные сценарии для криптоактивов. Но глобальную
универсальность и инфраструктурный охват карт они пока не
воспроизводят. Поэтому новые методы обычно входят в рынок рядом с
картами, там, где их частное преимущество
перевешивает общее карточное удобство.

\section{Структура способов оплаты в России и СНГ}~\label{sec:alt-mix}

Одна таблица с процентами создаёт иллюзию единой структуры способов
оплаты по всему региону. Устойчивыми остаются паттерны, а не цифры.

Россия сочетает массовый карточный контур и быстро растущие
альтернативы: Банк России прямо пишет, что карты остаются самым
популярным средством безналичной оплаты, а платежи по QR-коду
в 2024 году достигли 4{,}1~трлн рублей~\cite{cbr-results-2024-payments}. Казахстан показывает другой
профиль: национальная статистика отдельно выделяет крупный поток
безналичных QR-платежей наряду с карточными операциями~\cite{nbk-payment-cards-2025jan}. Узбекистан опирается на две
национальные карточные системы: HUMO и UZCARD~\cite{humo-faq,uzcard-about}. Армения сохраняет
собственный внутренний карточный контур через ArCa; оператор
системы прямо указывает, что банки-участники выпускают карты ArCa
наряду с Visa, Mastercard и другими международными
системами~\cite{arca-about}.

Из этого следует простой продуктовый вывод: СНГ нельзя считать единым
платёжным рынком. Локализуя продукт, нельзя механически переносить
российскую структуру способов оплаты на Казахстан или армянский
сценарий на Узбекистан без проверки локальной инфраструктуры, роли
QR и регуляторных требований.

\practicenote{%
При локализации в новой стране СНГ проверьте BIN-распределение
трафика, наличие массового локального A2A-метода и требования к
обязательному приёму локальных карт. Эти три фактора обычно сильнее
всего меняют экономику и конверсию.

}
